% ===== 原文 pp.292–306（PDF p.307–321），tr-03：第 6 章 §6.3 后段–§6.5 =====
% 块首为 §6.3 正文中段（§6.3 标题及更早内容属 c06a/c06b/c06c），故不写 §6.3 标题。

% [原文 p.292]
是 $\Bbbk[P]$ 在理想 $\mathfrak{m}$ 处的完备化。于是我们还可定义
\[
\hat{\mathbb{V}}:=\varprojlim\mathbb{V}_{\mathfrak{m}^k},
\]
它是 $\operatorname{Aut}(\widehat{\Bbbk[P]})$ 的一个 pro-幂零子群。

正如 §5.4.3 中那样，我们可以定义散射图的概念，定义基本相同。

\begin{definition}
给定 $r:P\to M$，并设 $I$ 是 $\mathfrak{m}$-准素理想。

(1) \emph{射线}或\emph{直线}是一对 $(\mathfrak{d},f_{\mathfrak{d}})$，使得
\begin{itemize}[leftmargin=1.6em]
\item $\mathfrak{d}\subseteq M_\R$ 由
\[
\mathfrak{d}=m_0'-\R_{\ge0}m_0
\]
给出（若 $\mathfrak{d}$ 是射线），或由
\[
\mathfrak{d}=m_0'-\R m_0
\]
给出（若 $\mathfrak{d}$ 是直线），其中 $m_0'\in M_\R$，$m_0\in M\setminus\{0\}$。集合 $\mathfrak{d}$ 称为该直线或射线的\emph{支撑}。若 $\mathfrak{d}$ 是射线，则 $m_0'$ 称为该射线的\emph{起点}，记作 $\operatorname{Init}(\mathfrak{d})$。
\item 对 $m_0\in M\setminus\{0\}$，令
\[
P_{m_0}:=\{m\in P\mid r(m)=Cm_0\text{ 对某个有理数 }C>0\}.
\]
则 $f_{\mathfrak{d}}\in\widehat{\Bbbk[P]}$ 满足
\[
f_{\mathfrak{d}}=1+\sum_{m\in P_{m_0}}c_mz^m.
\]
\item $f_{\mathfrak{d}}\equiv1\bmod\mathfrak{m}$。
\end{itemize}

(2) $\Bbbk[P]/I$ 上的一个\emph{散射图} $\mathfrak{D}$ 是有限条直线与射线之集，使得对每个 $(\mathfrak{d},f_{\mathfrak{d}})\in\mathfrak{D}$ 有 $f_{\mathfrak{d}}\in \Bbbk[P]$。$\widehat{\Bbbk[P]}$ 上的散射图 $\mathfrak{D}$ 是可数条直线与射线之集，使得对每个 $k\ge1$，只有有限条直线或射线 $(\mathfrak{d},f_{\mathfrak{d}})\in\mathfrak{D}$ 满足 $f_{\mathfrak{d}}\not\equiv1\bmod\mathfrak{m}^k$。
\end{definition}

% [Tutor Note: 定义 6.37 与第 5 章散射图定义的关键差别在“$\widehat{\Bbbk[P]}$ 上的可数性”条款：对每个 $k\ge1$ 只有有限条模 $\mathfrak{m}^k$ 非平凡的射线，故 $\theta_{\gamma,\mathfrak{D}}$ 的每次模 $\mathfrak{m}^k$ 计算都只涉及有限次复合（这正是取极限 $\theta_{\gamma,\mathfrak{D}}\in\hat{\mathbb{V}}$ 的合法性依据）。]

于是我们有与 §5.4.3 中定义的集合 $\operatorname{Supp}(\mathfrak{D})$、$\operatorname{Sing}(\mathfrak{D})$ 相同的概念，并且给定光滑浸入 $\gamma:[0,1]\to M_\R\setminus\operatorname{Sing}(\mathfrak{D})$，若固定一个 $\mathfrak{m}$-准素理想 $I$，则完全如 §5.4.3 那样得到自同构
\[
\theta_{\gamma,\mathfrak{D}}\in\mathbb{V}_I.
\]
确实，需要验证在穿过射线 $\mathfrak{d}_i$ 时出现的每个自同构 $\theta_{\gamma,\mathfrak{d}_i}$ 有意义；它的定义为 $\theta_{\gamma,\mathfrak{d}_i}(z^m)=z^mf_{\mathfrak{d}_i}^{\langle n,r(m)\rangle}$，其中 $n\in N$ 是通常的那个。这是有意义的，因为 $f_{\mathfrak{d}_i}\equiv1\bmod\mathfrak{m}$，故 $f_{\mathfrak{d}_i}$ 在 $\Bbbk[P]/I$ 中可逆。于是如前，我们得到自同构 $\theta_{\gamma,\mathfrak{D}}$，它只依赖于路径 $\gamma$ 在 $M_\R\setminus\operatorname{Sing}(\mathfrak{D})$ 中的同伦型。

此外，若 $\mathfrak{D}$ 是 $\widehat{\Bbbk[P]}$ 上的散射图，则 $\theta_{\gamma,\mathfrak{D}}$ 类似地有意义：对任意 $k>0$，计算 $\theta_{\gamma,\mathfrak{D}}$ 模 $\mathfrak{m}^k$ 是有意义的，因为这只涉及有限次复合。然后取极限便得到
\[
\theta_{\gamma,\mathfrak{D}}\in\hat{\mathbb{V}}.
\]
在 §5.4.3 中，我们由所有 Maslov 指标为零的圆盘之集构造了一个显式的散射图 $\mathfrak{D}$。我们发现，绕 $\mathfrak{D}$ 的某些奇点转圈的回路所对应的自同构 $\theta_{\gamma,\mathfrak{D}}$ 是恒等映射。在本节中，我们

% [原文 p.293]
将按恰好相反的方式来处理：先说明如何利用 Kontsevich 与 Soibelman 的论证 [70]，通过添加射线总能达到该条件。

\begin{theorem}
固定一个 $\mathfrak{m}$-准素理想 $I$。设 $\mathfrak{D}$ 是 $\Bbbk[P]/I$ 上的散射图。则存在包含 $\mathfrak{D}$ 的散射图 $S_I(\mathfrak{D})$，使得 $S_I(\mathfrak{D})\setminus\mathfrak{D}$ 只由射线组成，并且对任何使 $\theta_{\gamma,S_I(\mathfrak{D})}$ 有定义的闭回路 $\gamma$，有 $\theta_{\gamma,S_I(\mathfrak{D})}=\operatorname{Id}\in\mathbb{V}_I$。

类似地，若 $\mathfrak{D}$ 是 $\widehat{\Bbbk[P]}$ 上的散射图，则存在包含 $\mathfrak{D}$ 的散射图 $S(\mathfrak{D})$，使得 $S(\mathfrak{D})\setminus\mathfrak{D}$ 只由射线组成，并且对任何使 $\theta_{\gamma,S(\mathfrak{D})}$ 有定义的闭回路 $\gamma$，有 $\theta_{\gamma,S(\mathfrak{D})}=\operatorname{Id}\in\hat{\mathbb{V}}$。
\end{theorem}

\begin{proof}[证明]
我们对 $k$ 归纳，证明存在 $\mathfrak{D}_k$ 使得对任何使 $\theta_{\gamma,\mathfrak{D}_k}$ 有定义的闭回路 $\gamma$ 有
\[
\theta_{\gamma,\mathfrak{D}_k}\equiv\operatorname{Id}\bmod\mathfrak{m}^{k+1}.
\]
在 $\Bbbk[P]/I$ 上工作的情形，由于 $I$ 是 $\mathfrak{m}$-准素的，存在某个 $k$ 使 $\mathfrak{m}^{k+1}\subseteq I$，于是取 $S_I(\mathfrak{D})=\mathfrak{D}_k$ 即可。否则我们无限进行下去，取 $S(\mathfrak{D})$ 为诸 $\mathfrak{D}_k$ 的（非不交）并。

取 $\mathfrak{D}_0=\mathfrak{D}$。为从 $\mathfrak{D}_{k-1}$ 得到 $\mathfrak{D}_k$，我们如下进行。令 $\mathfrak{D}_{k-1}'$ 由 $\mathfrak{D}_{k-1}$ 中满足 $f_{\mathfrak{d}}\not\equiv1\bmod\mathfrak{m}^{k+1}$ 的那些射线与直线 $\mathfrak{d}$ 组成。注意 $\operatorname{Sing}(\mathfrak{D}_{k-1}')$ 是有限集。设 $p\in\operatorname{Sing}(\mathfrak{D}_{k-1}')$。令 $\gamma_p$ 为绕 $p$ 的闭简单回路，取足够小使得其内部不含 $\operatorname{Sing}(\mathfrak{D}_{k-1}')$ 的其他点。显然
\[
\theta_{\gamma_p,\mathfrak{D}_{k-1}}\equiv\theta_{\gamma_p,\mathfrak{D}_{k-1}'}\bmod\mathfrak{m}^{k+1}.
\]
由归纳假设，我们可以唯一地写成
\[
\theta_{\gamma_p,\mathfrak{D}_{k-1}'}\equiv\exp\Big(\sum_{i=1}^sc_iz^{m_i}\partial_{n_i}\Big)\bmod\mathfrak{m}^{k+1},
\]
其中 $m_i\in\mathfrak{m}^k$，$r(m_i)\ne0$，$n_i\in r(m_i)^\perp$ 本原，$c_i\in k$。令
\[
\mathfrak{D}[p]=\{(p-\R_{\ge0}r(m_i),1\pm c_iz^{m_i})\mid i=1,\dots,s\}.
\]
每条射线中的符号选取使得它对 $\theta_{\gamma_p,\mathfrak{D}[p]}$ 的贡献模 $\mathfrak{m}^{k+1}$ 为 $\exp(-c_iz^{m_i}\partial_{n_i})$。由于对任何对数导子 $\xi\in\mathfrak{v}$ 有 $[c_iz^{m_i}\partial_{n_i},\xi]\equiv0\bmod\mathfrak{m}^{k+1}$，$\mathfrak{D}[p]$ 中任一条射线对应的自同构模 $\mathfrak{m}^{k+1}$ 与 $\mathfrak{D}_{k-1}$ 中任一条射线对应的自同构交换。于是我们有
\[
\theta_{\gamma_p,\mathfrak{D}_{k-1}\cup\mathfrak{D}[p]}=\operatorname{Id}\bmod\mathfrak{m}^{k+1}.
\]
此外
\[
\mathfrak{D}_k=\mathfrak{D}_{k-1}\cup\bigcup_p\mathfrak{D}[p]
\]
现在具有所需性质。确实，设 $p\in\operatorname{Sing}(\mathfrak{D}_k)$。对 $\theta_{\gamma_p,\mathfrak{D}_k}$ 模 $\mathfrak{m}^{k+1}$ 有贡献、且不在 $\mathfrak{D}_{k-1}\cup\mathfrak{D}[p]$ 中的射线只可能属于 $\bigcup_{p'\ne p}\mathfrak{D}[p']$，而这些射线贡献两次，但对应的自同构互逆而相消。
\end{proof}

\begin{remark}
上述引理的证明中描述的过程非常简单，容易在计算机上实现。也容易看出结果是本质唯一的。唯一的非唯一性来源于可能存在若干条有相同支撑的不同射线或直线。给定射线或直线 $(\mathfrak{d}_1,f_{\mathfrak{d}_1}),\dots,(\mathfrak{d}_n,f_{\mathfrak{d}_n})$ 满足 $\operatorname{Supp}(\mathfrak{d}_1)=\cdots=\operatorname{Supp}(\mathfrak{d}_n)$，我们可以把这些射线或

% [原文 p.294]
直线换成 $(\mathfrak{d},\prod_{i=1}^nf_{\mathfrak{d}_i})$，其中 $d=\operatorname{Supp}(\mathfrak{d}_i)$。若对散射图 $\mathfrak{D}$ 中每一组这样的射线或直线施加该过程，便得到一个新图 $\mathfrak{D}'$，其中没有两个元素有相同的支撑，然而对任何使 $\theta_{\gamma,\mathfrak{D}}$ 有定义的路径，自同构 $\theta_{\gamma,\mathfrak{D}}$ 与 $\theta_{\gamma,\mathfrak{D}'}$ 相同。
\end{remark}

这便引出如下定义：

\begin{definition}
两个散射图 $\mathfrak{D}$ 与 $\mathfrak{D}'$ 在 $\Bbbk[P]/I$ 上（或在 $\widehat{\Bbbk[P]}$ 上）\emph{等价}，若对任何使 $\theta_{\gamma,\mathfrak{D}}$ 与 $\theta_{\gamma,\mathfrak{D}'}$ 都有定义的路径 $\gamma$，在 $\mathbb{V}_I$（或 $\hat{\mathbb{V}}$）中有 $\theta_{\gamma,\mathfrak{D}}=\theta_{\gamma,\mathfrak{D}'}$。
\end{definition}

正确的唯一性结果则是：$S_I(\mathfrak{D})$（或 $S(\mathfrak{D})$）在同构意义下唯一。

\begin{example}
$S(\mathfrak{D})$ 的结构可能非常复杂，即使 $\mathfrak{D}$ 的选取非常简单。我们在此讨论三个例子，读者可参阅 [44] 了解详情。

这些例子中我们取 $P=M\oplus N$，记 $x=z^{(1,0,0)}$，$y=z^{(0,1,0)}$，$t=z^{(0,0,1)}$，于是 $\Bbbk[P]= \Bbbk[x^{\pm1},y^{\pm1},t]$，$\mathfrak{m}=(t)$，且 $\widehat{\Bbbk[P]}= \Bbbk[x^{\pm1},y^{\pm1}][\![t]\!]$。

考虑散射图
\[
\mathfrak{D}=\{(\R_{(1,0)},(1+tx^{-1})^{\ell_1}),(\R_{(0,1)},(1+ty^{-1})^{\ell_2})\},
\]
其中 $\ell_1,\ell_2>0$ 任意。若 $\ell_1=\ell_2=1$，则容易手算验证
\[
S(\mathfrak{D})\setminus\mathfrak{D}=\{(\R_{(1,1)},1+t^2x^{-1}y^{-1})\}.
\]
确实，这与例 6.29 中作的计算相同，只需多一条射线。但若 $\ell_1=\ell_2=2$，就已经需要无穷多条射线：
\begin{align*}
S(\mathfrak{D})\setminus\mathfrak{D}={}&\{(\R_{(n+1,n)},(1+t^{2n+1}x^{-(n+1)}y^{-n})^2)\mid n\in\Z,\ n\ge1\}\\
&\cup\{(\R_{(n,n+1)},(1+t^{2n+1}x^{-n}y^{-(n+1)})^2)\mid n\in\Z,\ n\ge1\}\\
&\cup\{(\R_{(1,1)},(1-t^2x^{-1}y^{-1})^{-4})\}.
\end{align*}
对 $\ell_1=\ell_2=3$，情形要复杂得多。Kontsevich 首先注意到某种对称性：
\[
(\R_{\ge0}(m_1,m_2),f(t^{m_1+m_2}x^{-m_1}y^{-m_2}))\in S(\mathfrak{D})
\]
当且仅当
\[
(\R_{\ge0}(3m_1-m_2,m_1),f(t^{4m_1-m_2}x^{m_2-3m_1}y^{-m_1}))\in S(\mathfrak{D}),
\]
只要 $m_1,m_2$ 与 $3m_1-m_2$ 都为正。此外，$S(\mathfrak{D})$ 含射线
\[
(\R_{\ge0}(3,1),(1+t^4x^{-3}y^{-1})^3)\quad\text{与}\quad(\R_{\ge0}(1,3),(1+t^4x^{-1}y^{-3})^3),
\]
从而由上述对称性，还有支撑为
\[
\R_{\ge0}(8,3),\ \R_{\ge0}(21,8),\ \dots
\]
与
\[
\R_{\ge0}(3,8),\ \R_{\ge0}(8,21),\ \dots
\]
的射线，它们收敛到斜率为 $(3\pm\sqrt5)/2$ 的射线，这两个斜率对应于线性变换
\[
\begin{pmatrix}3&-1\\ 1&0\end{pmatrix}
\]
的两个不同特征空间。注意由对称性，我们知道这些射线上所附的函数是什么。另一方面，在斜率为 $(3\pm\sqrt5)/2$ 的射线所张成的锥内部，每个有理斜率都会出现，而这些射线上所附的函数看起来非常复杂。唯一已知

% [原文 p.295]
的情形是斜率为 $1$ 的直线（以及由对称性得到的斜率）上所附的函数，Reineke 在 [95] 中证明了它是
\[
\Big(\sum_{n=0}^\infty\frac{1}{3n+1}\binom{4n}{n}t^{2n}x^{-n}y^{-n}\Big)^9.
\]
虽然这些散射图非常复杂，但注意：对任何固定的 $k$，模 $t^k$ 只有有限个 $f_{\mathfrak{d}}$ 满足 $f_{\mathfrak{d}}\not\equiv1$。
\end{example}

\subsubsection{计数解释}

在第 5 章中，我们看到了“绕奇点的回路 $\gamma$ 满足 $\theta_{\gamma,\mathfrak{D}}=\operatorname{Id}$”这一条件与“把 Maslov 指标为零的圆盘粘合起来产生新的 Maslov 指标为零的圆盘”之间的关系。这里考虑的散射图要复杂得多，但事实上仍有类似的解释，这归功于我本人、Pandharipande 与 Siebert，见 [45]。虽然我会在此给出精确陈述（因为我们需要该结果的一个方面），但我不会给出证明的任何细节。因此本节基本上是本章论证中唯一不完全自足的部分。

我们考虑如下情形。设 $\ell_1,\dots,\ell_p>0$ 为 $p$ 个正整数，并令 $\ell=\sum_{i=1}^p\ell_i$。置 $P=M\oplus N^\ell$。我们把 $\Bbbk[P]$ 中对应于 $N^\ell$ 的生成元的单项式记为 $t_{ij}$，$1\le i\le p$，$1\le j\le\ell_i$。设 $m_1,\dots,m_p\in M$ 是互不相同的本原元素。令 $\mathfrak{D}$ 为散射图
\begin{equation}
\mathfrak{D}=\{(\R_{m_i},\prod_{j=1}^{\ell_i}(1+t_{ij}z^{-m_i}))\mid1\le i\le p\}.
\tag{6.6}
\end{equation}
理想 $\mathfrak{m}$ 由变量 $t_{ij}$ 生成。我们可以研究 $\widehat{\Bbbk[P]}$ 上的无穷散射图 $S(\mathfrak{D})$，即我们可以在环 $ \Bbbk[M][\![\{t_{ij}\}]\!]$ 上工作。特别地，在模去任何 $\mathfrak{m}$-准素理想 $I$ 之后，这将给出 $S_I(\mathfrak{D})$。确实，给定 $S(\mathfrak{D})$，我们可以通过删去那些满足 $f_{\mathfrak{d}}\equiv1\bmod I$ 的射线 $(\mathfrak{d},f_{\mathfrak{d}})\in S(\mathfrak{D})$，并对每个满足 $f_{\mathfrak{d}}\not\equiv1\bmod I$ 的 $(\mathfrak{d},f_{\mathfrak{d}})\in S(\mathfrak{D})$ 去掉 $f_{\mathfrak{d}}$ 中属于 $I$ 的单项式，从而得到 $\Bbbk[P]/I$ 上的有限散射图；这样 $f_{\mathfrak{d}}$ 就是多项式而非幂级数。

我们希望解释 $S(\mathfrak{D})\setminus\mathfrak{D}$ 的元素。设 $S(\mathfrak{D})$ 选取得使没有两条射线有相同支撑，并设 $(\mathfrak{d},f_{\mathfrak{d}})\in S(\mathfrak{D})\setminus\mathfrak{D}$。更具体地说，我们将解释 $f_{\mathfrak{d}}$。于是设讨论中 $\mathfrak{d}$ 已固定。

设 $\Sigma_{\mathfrak{d}}$ 是 $M_\R$ 中定义非异环面曲面 $X_{\mathfrak{d}}$ 的完备扇，具有性质：在其一维射线的集合中包含
\[
\R_{\ge0}m_1,\dots,\R_{\ge0}m_p\text{ 以及 }\mathfrak{d}.
\]
（注意 $\mathfrak{d}$ 可能与其他射线重合。）$\Sigma_{\mathfrak{d}}$ 的精确选取最终将无关紧要。

设 $D_1,\dots,D_p$ 为 $X_{\mathfrak{d}}$ 中对应于射线 $\R_{\ge0}m_1,\dots,\R_{\ge0}m_p$ 的环面除子，并设 $D_{\mathrm{out}}$ 对应于射线 $\mathfrak{d}$（可能对某个 $i$ 有 $D_i=D_{\mathrm{out}}$）。选取一般点 $\{x_{i1},\dots,x_{i\ell_i}\}\subseteq D_i$。令
\[
\nu:\tilde X_{\mathfrak{d}}\to X_{\mathfrak{d}}
\]
为 $X_{\mathfrak{d}}$ 在点集 $\{x_{ij}\}$ 处的爆破。设 $\tilde D_1,\dots,\tilde D_p,\tilde D_{\mathrm{out}}$ 为相应环面除子的真变换，并设 $E_{ij}$ 为点 $x_{ij}$ 上的例外除子。

% [原文 p.296]
现在引入附加数据 $\mathbf{P}=(\mathbf{P}_1,\dots,\mathbf{P}_p)$，其中 $\mathbf{P}_i$ 表示一个由 $\ell_i$ 个非负整数构成的序列 $(p_{i1},\dots,p_{i\ell_i})$。我们将使用记号 $\mathbf{P}_i=p_{i1}+\cdots+p_{i\ell_i}$，并称 $\mathbf{P}_i$ 为一个\emph{有序分拆}。我们称 $p_{i1},\dots,p_{i\ell_i}$ 为 $\mathbf{P}_i$ 的\emph{部分}。定义
\[
|\mathbf{P}_i|=\sum_{j=1}^{\ell_i}p_{ij},
\]
并限制到满足
\begin{equation}
-\sum_{i=1}^p|\mathbf{P}_i|m_i=k_{\mathbf{P}}m_{\mathfrak{d}}
\tag{6.7}
\end{equation}
的那些 $\mathbf{P}$，其中 $m_{\mathfrak{d}}\in M$ 是 $\mathfrak{d}$ 的本原生成元，$k_{\mathbf{P}}$ 是正整数。

给定这些数据，考虑由如下要求确定的类 $\beta\in H_2(X_{\mathfrak{d}},\Z)$：若 $D$ 是 $X_{\mathfrak{d}}$ 的环面除子且 $D\notin\{D_1,\dots,D_p,D_{\mathrm{out}}\}$，则 $D\cdot\beta=0$；若 $D_{\mathrm{out}}\notin\{D_1,\dots,D_p\}$，则
\[
D_i\cdot\beta=|\mathbf{P}_i|,\qquad D_{\mathrm{out}}\cdot\beta=k_{\mathbf{P}};
\]
而若对某个 $j$ 有 $D_{\mathrm{out}}=D_j$，则
\[
D_i\cdot\beta=\begin{cases}|\mathbf{P}_i|,&i\ne j\\ |\mathbf{P}_i|+k_{\mathbf{P}},&i=j.\end{cases}
\]
这样的类存在，这由 (6.7) 与命题 4.2 推出。于是我们可以定义
\[
\beta_{\mathbf{P}}=\nu^*(\beta)-\sum_{i=1}^p\sum_{j=1}^{\ell_i}p_{ij}[E_{ij}]\in H_2(\tilde X_{\mathfrak{d}},\Z).
\]
粗略地说，这就是 $\tilde X_{\mathfrak{d}}$ 中一条曲线的真变换的类，该曲线以与 $D_i$ 横截的方式 $p_{ij}$ 次经过点 $x_{ij}$，并以计重数的方式与 $D_{\mathrm{out}}$ 交于 $k_{\mathbf{P}}$ 个点。最后，令 $X_{\mathfrak{d}}^o$ 由 $X_{\mathfrak{d}}$ 去掉所有零维环面轨道得到，并令 $\tilde X_{\mathfrak{d}}^o=\nu^{-1}(X_{\mathfrak{d}}^o)$，$\tilde D_{\mathrm{out}}^o=\tilde D_{\mathrm{out}}\cap\tilde X_{\mathfrak{d}}^o$。

我们现在（略不正式地）考虑表示类 $\beta_{\mathbf{P}}$ 的亏格零相对稳定映射的模空间
\[
\overline{\mathfrak{M}}(\tilde X_{\mathfrak{d}}^o/\tilde D_{\mathrm{out}}^o)\subseteq\overline{\mathfrak{M}}(\tilde X_{\mathfrak{d}}/\tilde D_{\mathrm{out}}),
\]
其中我们施加与 $\tilde D_{\mathrm{out}}$ 的极大相切性。粗略地说，$\overline{\mathfrak{M}}(\tilde X_{\mathfrak{d}}/\tilde D_{\mathrm{out}})$ 是稳定映射 $f:C\to\tilde X_{\mathfrak{d}}$ 的模空间，这些映射满足：
\begin{enumerate}[label=(\arabic*)]
\item $C$ 亏格为零。
\item $f$ 表示类 $\beta_{\mathbf{P}}$。
\item 存在唯一的点 $p\in C$ 使 $f(p)\in\tilde D_{\mathrm{out}}$，且若 $x$ 是 $\tilde D_{\mathrm{out}}$ 在 $f(p)$ 附近的局部方程，则 $f^*(x)$ 有 $k_{\mathbf{P}}$ 阶零点。
\end{enumerate}
一般地，这并不完全准确，因为这样的映射空间不是紧的。为正确给出该模空间，必须允许不仅对曲线 $C$、也对空间 $\tilde X_{\mathfrak{d}}$ 发生冒泡现象。这在代数几何语境中由 Jun Li 在 [72] 与 [73] 中完成。参见 [35] 对该构造的综述。不过不熟悉这一概念的读者不必过分担心，因为此处不需要该构造的细节。

于是我们定义 $\overline{\mathfrak{M}}(\tilde X_{\mathfrak{d}}^o/\tilde D_{\mathrm{out}}^o)$ 为由那些像包含在 $\tilde X_{\mathfrak{d}}^o$（而非 $\tilde X_{\mathfrak{d}}$）中的相对稳定映射构成的开子空间。虽然 $\overline{\mathfrak{M}}(\tilde X_{\mathfrak{d}}/\tilde D_{\mathrm{out}})$

% [原文 p.297]
因目标空间 $\tilde X_{\mathfrak{d}}$ 紧而紧，但事实上开子集 $\overline{\mathfrak{M}}(\tilde X_{\mathfrak{d}}^o/\tilde D_{\mathrm{out}}^o)$ 也是紧的：这样的曲线不可能收敛到经过某个被删去点的曲线。这在 [45] 的命题 5.1 中证明。此外，空间 $\overline{\mathfrak{M}}(\tilde X_{\mathfrak{d}}^o/\tilde D_{\mathrm{out}}^o)$（一般是 Deligne–Mumford 叠）有期望维数零与虚拟基本类。于是我们定义数
\[
N_{\mathbf{P}}:=\int_{[\overline{\mathfrak{M}}(\tilde X_{\mathfrak{d}}^o/\tilde D_{\mathrm{out}}^o)]^{\mathrm{vir}}}1.
\]
不正式地说，这就是类 $\beta_{\mathbf{P}}$ 中与 $\tilde D_{\mathrm{out}}$ 极大相切的有理曲线之数。当然，这些数可能更微妙，涉及 Gromov–Witten 理论中通常出现的虚拟计数：多重覆盖与叠现象可能产生分数。

有了这些定义，我们现在可以陈述 [45] 的主要定理：

\begin{theorem}
\[
\log f_{\mathfrak{d}}=\sum_{\mathbf{P}}k_{\mathbf{P}}N_{\mathbf{P}}t_{\mathbf{P}}z^{-k_{\mathbf{P}}m_{\mathfrak{d}}},
\]
其中求和遍历所有满足 (6.7) 的 $\mathbf{P}$，且 $t_{\mathbf{P}}$ 表示单项式 $\prod_{ij}t_{ij}^{p_{ij}}$。
\end{theorem}

% [Tutor Note: 定理 6.42 是 §6.3.2 唯一的非自足输入：它把散射图射线上所附函数 $f_{\mathfrak{d}}$ 的对数展开系数解释为相对 Gromov–Witten 不变量 $N_{\mathbf{P}}$。取 $\log$ 而不是 $\exp$ 的原因可见例 6.43(1)：$\log(1+t_{11}t_{21}x^{-1}y^{-1})$ 的系数是 $(-1)^{k+1}/k^2$，即直线的多重覆盖以 $1/k^2$ 型分母记录。]

\begin{examples}
(1) 我们考虑 $p=2$、$m_1=(1,0)$、$m_2=(0,1)$ 的情形；这本质上就是例 6.41 中考虑的情形，只是把变量 $t$ 换成变量 $t_{11},\dots,t_{1\ell_1},t_{21},\dots,t_{2\ell_2}$。首先，当 $\ell_1=\ell_2=1$ 时，我们有一条输出射线 $(\R_{\ge0}(1,1),1+t_{11}t_{21}x^{-1}y^{-1})$。（要得到例 6.41 中原来的 $1+t^2x^{-1}y^{-1}$，只需代换 $t=t_{11}=t_{12}$。）射线 $\R_{\ge0}(-1,0)$、$\R_{\ge0}(0,-1)$ 与 $\R_{\ge0}(1,1)$ 是定义 $\PP^2$ 的扇中的一维锥，而三个环面除子 $D_1,D_2,D_{\mathrm{out}}$ 就是三条坐标轴。在 $D_1$ 上固定一点、在 $D_2$ 上固定一点，恰有一条直线经过这两点；取 $\mathbf{P}=(1,1)$，我们便看到 $N_{\mathbf{P}}=1$。另一方面注意
\[
\log(1+t_{11}t_{21}x^{-1}y^{-1})=\sum_{k=1}^\infty\frac{(-1)^{k+1}}{k^2}kt_{11}^kt_{21}^kx^{-k}y^{-k}.
\]
于是我们看到 $N_{(k,k)}=(-1)^{k+1}/k^2$。对 $k=1$，这就是上面计数的那条单条直线；其余贡献来自该直线的多重覆盖，这些覆盖在直线与 $D_{\mathrm{out}}$ 的交点处完全分歧。

(2) 其次考虑 $\ell_1=\ell_2=2$ 的情形，并考虑例 6.41 中描述的 $S(\mathfrak{D})$ 里的输出射线 $(\R_{\ge0}(1,1),(1-t^2x^{-1}y^{-1})^{-4})$。注意
\[
\log(1-t^2x^{-1}y^{-1})^{-4}=4\sum_{k=1}^\infty k\cdot\frac{1}{k^2}t^{2k}x^{-k}y^{-k}.
\]
由于我们把四个变量 $t_{ij}$（$1\le i,j\le2$）都代换为 $t$，每个系数代表对多种 $\mathbf{P}$ 的求和。特别地，$N_{(1+0,1+0)}+N_{(0+1,1+0)}+N_{(1+0,0+1)}+N_{(0+1,0+1)}=4$ 是 $t^2x^{-1}y^{-1}$ 的系数。这是经过两个点的直线之数，其中一个点选在 $\{x_{11},x_{12}\}$ 中，另一个选在 $\{x_{21},x_{22}\}$ 中。$t^4x^{-2}y^{-2}$ 的系数是 $1$，即 $N_{(1+1,1+1)}+N_{(2+0,2+0)}+\cdots+N_{(0+2,0+2)}$。这里形如 $N_{(2+0,1+1)}$ 的不变量容易看出为零，而形如 $N_{(2+0,2+0)}$ 的只给出直线的多重覆盖贡献，对四条直线各为 $-1/4$。于是 $N_{(1+1,1+1)}=2$，这解释了经过四个点 $\{x_{ij}\}$ 且与 $D_{\mathrm{out}}$ 相切的二次曲线之数。

(3) 最后看 $\ell_1=\ell_2=3$ 的情形，仍看斜率为 $1$ 的射线，其上附有函数 $f_{\mathfrak{d}}$。则由例 6.41 中给出的显式公式可算得

% [原文 p.298]
\[
\log f_{\mathfrak{d}}=9t^2x^{-1}y^{-1}+2\cdot\frac{63}{4}t^4x^{-2}y^{-2}+3\cdot55t^6x^{-3}y^{-3}+\cdots.
\]
账目如下。第一个系数 $9$ 计数经过 $\{x_{11},x_{12},x_{13}\}$ 中一点与 $\{x_{21},x_{22},x_{32}\}$ 中一点的直线数。第二个系数 $63/4$ 计入这些直线的二重覆盖，共 $-9/4$，以及经过这两组中各两点且与 $D_{\mathrm{out}}$ 相切的二次曲线之数，共 $3\times3\times2=18$。注意 $18-9/4=63/4$。三次曲线的账目（可通过不用单个变量 $t$ 而用变量 $t_{ij}$ 计算散射图来验证）如下。有 $N_{(1+1+1,1+1+1)}=18$，这是经过全部六个点 $\{x_{ij}\}$ 且与 $D_{\mathrm{out}}$ 三重相切的结点三次曲线之数。数 $N_{(2+1+0,1+1+1)}=3$：这是与 $D_{\mathrm{out}}$ 三重相切、其结点与 $x_{11}$ 重合、且经过除 $x_{13}$ 之外全部六个点的结点三次曲线之数。注意这种类型的 $\mathbf{P}$ 有 $2\times3\times2=12$ 种选取，故来自真正三次曲线的计数为 $18+36=54$。另一方面，我们还有 $9$ 条直线的三重覆盖，各贡献 $1/9$，故总数为 $54+1=55$。
\end{examples}

事实上，[45] 对更一般的散射图给出了公式。假设对每个 $1\le i\le p$ 给定正整数 $p_i$，并对每个 $1\le j\le p_i$ 给定非负整数 $\ell_{ij}$。考虑
\begin{equation}
\mathfrak{D}=\Big\{\Big(\R_{m_i},\prod_{j=1}^{p_i}\prod_{k=1}^{\ell_{ij}}(1+t_{ijk}z^{-jm_i})\Big)\Big|1\le i\le p\Big\}.
\tag{6.8}
\end{equation}
如前，$m_1,\dots,m_p\in M$ 是互不相同的本原元素。我们再次考虑 $ \Bbbk[M][\![\{t_{ijk}\}]\!]$ 上的 $S(\mathfrak{D})$，并解释射线 $(\mathfrak{d},f_{\mathfrak{d}})\in S(\mathfrak{D})$。

我们定义\emph{分次分拆}为一个有序分拆的有限序列 $\mathbf{G}=(\mathbf{P}_1,\dots,\mathbf{P}_d)$，其中 $\mathbf{P}_i$ 的每个部分都可被 $i$ 整除。我们写 $|\mathbf{G}|=\sum_{i=1}^d|\mathbf{P}_d|$。设 $\mathbf{G}=(\mathbf{G}_1,\dots,\mathbf{G}_p)$ 为一个分次分拆的 $p$-元组，其中 $\mathbf{G}_i$ 的第 $j$ 块 $\mathbf{P}_{ij}$ 形如 $p_{ij1}+\cdots+p_{ij\ell_{ij}}$。限制到满足
\begin{equation}
-\sum_{i=1}^p|\mathbf{G}_i|m_i=k_{\mathbf{G}}m_{\mathfrak{d}}
\tag{6.9}
\end{equation}
（对某个正整数 $k_{\mathbf{G}}$）的那些 $\mathbf{G}$。给定这些数据，考虑由如下要求确定的类 $\beta\in H_2(X_{\mathfrak{d}},\Z)$：若 $D\notin\{D_1,\dots,D_p,D_{\mathrm{out}}\}$，则 $D\cdot\beta=0$；若 $D_{\mathrm{out}}\notin\{D_1,\dots,D_p\}$，则
\[
D_i\cdot\beta=|\mathbf{G}_i|,\qquad D_{\mathrm{out}}\cdot\beta=k_{\mathbf{G}};
\]
而若对某个 $j$ 有 $D_{\mathrm{out}}=D_j$，则
\[
D_i\cdot\beta=\begin{cases}|\mathbf{G}_i|,&i\ne j\\ |\mathbf{G}_i|+k_{\mathbf{G}},&i=j.\end{cases}
\]
取一般点 $x_{ijk}\in D_i$，$1\le i\le p$，$1\le j\le p_i$，$1\le k\le\ell_{ij}$。下面我们将比之前略欠精确。定义 $N_{\mathbf{G}}$ 为如下有理曲线的虚拟数：它与 $D_{\mathrm{out}}$ 恰好在一个点处相切 $k_{\mathbf{G}}$ 阶，并且对每个点 $x_{ijk}$，曲线有 $p_{ijk}/j$ 个分支与 $D_i$ 相切 $j$ 阶。

% [原文 p.299]
后一个关联条件与之前一样，通过爆破每个点 $x_{ijk}$ 得到，但当 $j>1$ 时对点 $x_{ijk}$ 作的是 orbifold 爆破而非普通爆破，详见 [45] §5.5。于是我们有 [45] §5.7 中的类似定理：

\begin{theorem}
\[
\log f_{\mathfrak{d}}=\sum_{\mathbf{G}}k_{\mathbf{G}}N_{\mathbf{G}}t_{\mathbf{G}}z^{-k_{\mathbf{G}}m_{\mathfrak{d}}},
\]
其中求和遍历所有满足 (6.9) 的 $\mathbf{G}$，且 $t_{\mathbf{G}}$ 表示单项式 $\prod t_{ijk}^{p_{ijk}/j}$。
\end{theorem}

\begin{remark}
$\overline{\mathfrak{M}}(\tilde X_{\mathfrak{d}}^o/\tilde D_{\mathrm{out}}^o)$ 中的相对稳定映射可能是相当复杂的对象，涉及映射 $f:C\to\hat X_{\mathfrak{d}}^o$，其中 $\hat X_{\mathfrak{d}}^o$ 是可约概形。然而该可约概形带有映射 $\hat X_{\mathfrak{d}}^o\to\tilde X_{\mathfrak{d}}^o$，把 $f$ 与该映射复合，便得到映射 $\tilde f:C\to\tilde X_{\mathfrak{d}}$。由 [44] 的引理 4.1 与 4.2 的论证可知，$\tilde f$ 的像在爆破映射 $\nu:\tilde X_{\mathfrak{d}}\to X_{\mathfrak{d}}$ 之下只与环面边界的真变换交于 $\tilde D_{\mathrm{out}}^o$。这里 $\nu$ 在定理 6.42 所考虑的情形是 $X_{\mathfrak{d}}$ 在点 $x_{ij}$ 处的普通爆破，在定理 6.44 所考虑的情形则是点 $x_{ijk}$ 处的加权爆破。
\end{remark}

我们还可把 $\tilde f$ 与炸开映射 $\nu$ 复合，此时得到映射 $\bar f:C\to X_{\mathfrak{d}}$。注意 $\tilde f$ 与 $\bar f$ 都不必稳定。定义域 $C$ 可能有许多不同的分量，$\bar f$ 在其中一些分量上可能是常值。然而，由上述观察：$\tilde f$ 的像与 $\tilde D_i$ 不交，以及 $\tilde f(C)$ 是类 $\beta_{\mathbf{P}}$ 中的除子（或在更一般情形是类似定义的类 $\beta_{\mathbf{G}}$），可知除子 $\bar f_*(C)$ 在 $X_{\mathfrak{d}}$ 上与 $D_i$ 在点 $x_{ij}$ 处的相交重数恰为 $p_{ij}$（或在点 $x_{ijk}$ 处恰为 $p_{ijk}$）。此外，虽然 $\bar f_*(C)$ 可能有多个不可约分量，但存在点 $q\in D_{\mathrm{out}}$ 使
\[
\bar f_*(C)\cap\partial X_{\mathfrak{d}}=\{q\}\cup\{x_{ij}\mid1\le i\le p,\ 1\le j\le\ell_i\}
\]
（或
\[
\bar f_*(C)\cap\partial X_{\mathfrak{d}}=\{q\}\cup\{x_{ijk}\mid1\le i\le p,\ 1\le j\le p_i,\ 1\le k\le\ell_{ij}\}).
\]
点 $q$ 还受一个约束，我们在更一般的情形中描述它。确实，$\bar f_*(C)$ 是表示 $\beta\in H_2(X_{\mathfrak{d}},\Z)\cong H_2(X_{\mathfrak{d}},\Z)\cong\operatorname{Pic}X_{\mathfrak{d}}$ 的除子（因 $X_{\mathfrak{d}}$ 是二维有理的），而把该元素限制到 $\partial X_{\mathfrak{d}}$ 给出 $\operatorname{Pic}(\partial X_{\mathfrak{d}})$ 的一个确定元素。由于 $\bar f_*(C)$ 在由 $\beta$ 定义的线性系中，$\bar f_*(C)|_{\partial X_{\mathfrak{d}}}$ 在由 $\beta|_{\partial X_{\mathfrak{d}}}$ 给出的线性等价类中。注意
\[
\bar f_*(C)|_{\partial X_{\mathfrak{d}}}=\sum_{i,j,k}p_{ijk}x_{ijk}+k_{\mathbf{G}}q.
\]
由于 $\partial X_{\mathfrak{d}}$ 是（奇）椭圆曲线，对固定的 $x_{ijk}$ 选取，只有有限多个 $q$ 使该除子属于由 $\beta|_{\partial X_{\mathfrak{d}}}$ 给出的线性等价类。

\begin{remark}
虽然散射图 (6.8) 看起来仍相当特殊，但事实上它可用来给出定理 6.38 对几乎任何散射图的结论。我们在如下情形中解释这一过程：给定某个 $P$、照常的 $r:P\to M$，以及散射图
\[
\mathfrak{D}=\{(\R_{r(m_i)},f_i)\},
\]
% [原文 p.300]
其中 $f_i=1+c_iz^{-m_i}$，$c_i\in k$，$m_i\in P$ 且 $r(m_i)\ne0$。我们希望描述 $\mathfrak{m}$-准素理想 $I$ 的 $S_I(\mathfrak{D})$。

首先，$\mathfrak{D}$ 可能有一些重合的直线，故我们如下用一个等价的散射图替换 $\mathfrak{D}$（并重新命名诸 $m_i$）。存在互不相同的本原向量 $\bar m_1,\dots,\bar m_p\in M$、对每个 $1\le i\le p$ 的正整数 $p_i$，以及对每个 $1\le i\le p$、$1\le j\le p_i$ 的正整数 $\ell_{ij}$，使得
\[
\mathfrak{D}=\Big\{\Big(\R_{\bar m_i},\prod_{j=1}^{p_i}\prod_{k=1}^{\ell_{ij}}(1+c_{ijk}z^{-m_{ijk}})\Big)\Big|1\le i\le p\Big\},
\]
其中 $r(m_{ijk})$ 与 $\bar m_i$ 正比，且 $r(m_{ijk})$ 的指数为 $j$。令 $\ell=\sum_{ij}\ell_{ij}$，并置
\[
\mathfrak{D}'=\Big\{\Big(\R_{\bar m_i},\prod_{j=1}^{p_i}\prod_{k=1}^{\ell_{ij}}(1+t_{ijk}z^{-r(m_{ijk})})\Big)\Big|1\le i\le p\Big\}
\]
为环 $ \Bbbk[M][\![\{t_{ijk}\}]\!]$ 上的散射图；这就是 (6.8) 中考虑的散射图。于是 $S(\mathfrak{D}')$ 的射线由定理 6.44 给出。

设 $Q\subseteq M\oplus N^\ell$ 为由
\[
\{(-r(m_{ijk}),e_{ijk})\}
\]
自由生成的子幺半群，其中 $e_{ijk}\in N^\ell$ 对应于 $t_{ijk}$。则 $\mathfrak{D}'$ 事实上定义在 $ \Bbbk[M\oplus N^\ell]$ 的子环 $ \Bbbk[Q]$ 上。定义环同态
\[
\phi: \Bbbk[Q]\to \Bbbk[P],\qquad t_{ijk}z^{-r(m_{ijk})}\mapsto c_{ijk}z^{-m_{ijk}}.
\]
令 $I'=\phi^{-1}(I)$。我们得到 $ \Bbbk[Q]/I'$ 上的散射图 $S_{I'}(\mathfrak{D}')$，由唯一性可把它取为把 $S(\mathfrak{D}')$ 模 $I'$ 约化所得。

把 $\phi$ 作用于 $S_{I'}(\mathfrak{D}')$ 每个元素上所附的函数，我们得到定义在 $\Bbbk[P]$ 上的散射图 $\phi(S_{I'}(\mathfrak{D}'))$。此外，对绕原点的回路 $\gamma$，$\theta_{\gamma,\phi(S_{I'}(\mathfrak{D}'))}\equiv\operatorname{Id}\bmod I$，因为 $\theta_{\gamma,S_{I'}(\mathfrak{D}')}$ 在 $ \Bbbk[Q]/I'$ 中是恒等映射。于是由唯一性，事实上 $\phi(S_{I'}(\mathfrak{D}'))$ 必与 $S_I(\mathfrak{D})$ 等价。在这个意义上，$S(\mathfrak{D}')$ 是万有的散射图。

注意定理 6.38 证明中描述的那个特定过程所产生的射线上所附函数是二项式 $1+c_{\mathfrak{d}}z^{m_{\mathfrak{d}}}$。我们在此也可做到这一点：给定一条射线 $(\mathfrak{d},f_{\mathfrak{d}})\in\phi(S_{I'}(\mathfrak{D}'))$，我们可以用如下机制把 $f_{\mathfrak{d}}$ 在 $\Bbbk[P]/I$ 中唯一分解为这类因子之积。假设已分解
\[
f_{\mathfrak{d}}\equiv\prod_{i=1}^d(1+c_iz^{p_i})\bmod\mathfrak{m}^k.
\]
则模 $\mathfrak{m}^{k+1}$，
\[
f_{\mathfrak{d}}-\prod_{i=1}^d(1+c_iz^{p_i})\equiv\sum_jd_jz^{q_j},
\]
其中 $z^{q_j}\in\mathfrak{m}^k\setminus\mathfrak{m}^{k+1}$。于是模 $\mathfrak{m}^{k+1}$ 可写
\[
f_{\mathfrak{d}}\equiv\Big(\prod_{i=1}^d(1+c_iz^{p_i})\Big)\Big(\prod_j(1+d_jz^{q_j})\Big).
\]

% [原文 p.301]
我们如此继续，直到达到某个使 $\mathfrak{m}^k\subseteq I$ 的 $k$。于是我们可以把单条射线 $(\mathfrak{d},f_{\mathfrak{d}})$ 换成射线组 $\{(\mathfrak{d},1+c_iz^{p_i})\}$。这个新散射图满足定理 6.38 的结论，事实上它正是那里所给算法产生的散射图。

\end{remark}

我们需要一个容易由定理 6.42 得到的技术性观察。

\begin{proposition}
设 $\mathfrak{D}$ 为 (6.8) 中给出的散射图。设 $J$ 是 $R= \Bbbk[\{t_{ijk}\mid(i,j,k)\ne(1,1,1)\}]$ 中的单项式理想，使得 $R/J$ 是 Artin 的。也把 $J$ 记为它在 $ \Bbbk[M]\otimes_ k \Bbbk[\{t_{ijk}\}]$ 中生成的理想，并令 $\tilde I=(t_{111}^e)+J$ 为后一环中的单项式理想，于是 $\tilde I$ 是 $\mathfrak{m}$-准素理想。设 $\tilde{\mathfrak{D}}_e$ 为按定理 6.38 证明中的算法得到的散射图 $S_{\tilde I}(\mathfrak{D})$，使得 $S_{\tilde I}(\mathfrak{D})$ 不含形如 $(\mathfrak{d},f_{\mathfrak{d}})$ 且满足 $f_{\mathfrak{d}}\equiv1\bmod\tilde I$ 的元素。注意按构造有 $\tilde{\mathfrak{D}}_e\subseteq\tilde{\mathfrak{D}}_{e+1}$。则序列 $\tilde{\mathfrak{D}}_1,\tilde{\mathfrak{D}}_2,\dots$ 稳定。
\end{proposition}

\begin{proof}[证明]
考虑由如下分次分拆组 $\mathbf{G}=(\mathbf{G}_1,\dots,\mathbf{G}_p)$ 构成的集合 $\Gamma$：
\[
\prod_{\substack{i,j,k\\(i,j,k)\ne1}}t_{ijk}^{p_{ijk}/j}\notin J
\]
且对某个 $(i,j,k)\ne(1,1,1)$ 有 $p_{ijk}>0$。由于 $R/J$ 是 Artin 的，对 $p_{111}$ 的任一给定选取，只有有限个 $\mathbf{G}\in\Gamma$ 具有该 $p_{111}$；但 $\Gamma$ 本身是无限的，因为 $p_{111}$ 无界。特别地，若 $\mathbf{G}\in\Gamma$，则除 $p_{111}$ 外所有 $p_{ijk}$ 都被某个数 $N$ 界住。我们首先证明只有有限个 $\mathbf{G}\in\Gamma$ 满足 $N_{\mathbf{G}}\ne0$。

设 $\mathbf{G}\in\Gamma$ 且 $N_{\mathbf{G}}\ne0$。存在本原的 $m_{\mathfrak{d}}\in M$ 使 (6.9) 成立，从而对某个 $e$ 在散射图 $\tilde{\mathfrak{D}}_e$ 中有一条支撑为 $\mathfrak{d}=\R_{\ge0}m_{\mathfrak{d}}$ 的射线 $(\mathfrak{d},f_{\mathfrak{d}})$。照常取 $\Sigma_{\mathfrak{d}}$。我们总可以加细 $\Sigma_{\mathfrak{d}}$，故可假设 $\Sigma_{\mathfrak{d}}$ 同时含射线 $\R_{\ge0}(-m_1)$ 与 $\R_{\ge0}(m_1)$，于是投影 $M_\R\to M_\R/\R m_1$ 给出环面态射 $\pi:X_{\mathfrak{d}}\to\PP^1$。该态射有由对应于 $\R_{\ge0}m_1$ 的除子 $D_1$ 给出的截面。

由注 6.45，$N_{\mathbf{G}}\ne0$ 意味着存在映射 $\bar f:C\to X_{\mathfrak{d}}$，使除子 $\bar f_*(C)$ 与 $D_1$ 在 $x_{111}$ 处的相交重数为 $p_{111}$。此外，我们可以假设纤维 $\pi^{-1}(\pi(x_{111}))$ 不含在 $\bar f$ 的像中。确实，该纤维与 $\partial X_{\mathfrak{d}}$ 交于 $x_{111}$ 之外的另一点，而假设诸 $x_{ijk}$ 一般选取，我们可以假设该点不与任何 $x_{ijk}$ 重合。于是，在注 6.45 的记号下，我们可以称该点为 $q$。

现在由该注，$\partial X_{\mathfrak{d}}$ 上由 $\sum p_{ijk}x_{ijk}+k_{\mathbf{G}}q$ 给出的除子的除子类完全由 $\mathbf{G}$ 决定。然而存在某个 $(i,j,k)\ne(1,1,1)$ 使 $p_{ijk}\ne0$，故通过变动该 $x_{ijk}$ 而保持其余点不动，我们得到矛盾。于是对该 $x_{ijk}$ 的一般选取，$\bar f_*(C)$ 不能包含 $\pi^{-1}(\pi(x_{111}))$。

现在 $\bar f_*(C)$ 表示类 $\beta$，而 $\tilde f_*(C)$ 表示类 $\beta_{\mathbf{G}}=\nu^*(\beta)-p_{111}E_{111}-\cdots$。由于纤维 $\pi^{-1}(\pi(x_{111}))$ 的全变换由两个不可约分量组成，其中之一是 $E_{111}$，且 $E_{111}\cdot\beta_{\mathbf{G}}=p_{111}$，$\pi^{-1}(\pi(x_{111}))$ 与 $\bar f_*(C)$ 的相交重数必须至少为 $p_{111}$。由此我们得出：交数 $\beta\cdot F$（其中 $F$ 是 $\pi$ 的一条纤维的类）至少为 $p_{111}$。

% [原文 p.302]
其次注意，$\pi^{-1}(0)$ 与 $\pi^{-1}(\infty)$ 中至少有一个不含除子 $D_{\mathrm{out}}$；假设是 $\pi^{-1}(0)$。则 $\pi^{-1}(0)$ 在 $\tilde X_{\mathfrak{d}}$ 中的真变换与 $\tilde f(C)$ 不交。于是 $\beta\cdot\pi^{-1}(0)$ 完全由 $\mathbf{G}_2,\dots,\mathbf{G}_p$ 决定，因而有界。由于 $\pi^{-1}(0)$ 与 $F$ 线性等价，我们看到 $p_{111}$ 有界。该论证不依赖于 $\mathfrak{d}$，故 $p_{111}$ 与 $\mathfrak{d}$ 无关地有界。因此集合 $\{\mathbf{G}\in\Gamma\mid N_{\mathbf{G}}\ne0\}$ 有限。

特别地，由定理 6.44 与 (6.9) 可知，在 $\bigcup_e\tilde{\mathfrak{D}}_e$ 中至多出现有限多个不同的射线方向 $\mathfrak{d}$。

现在假设我们取 $\tilde{\mathfrak{D}}_e$ 使得任何给定支撑上至多有一条射线，于是对射线 $(\mathfrak{d},f_{\mathfrak{d}})\in\tilde{\mathfrak{D}}_e$，$\log f_{\mathfrak{d}}$ 由定理 6.44 给出。考虑在某 $\tilde{\mathfrak{D}}_e$ 中出现的射线 $(\mathfrak{d},f_{\mathfrak{d}})$。由于它是射线，任何系数都必须涉及某个 $(i,j,k)\ne1$ 的 $t_{ijk}$；否则它只能出现在计算 $S(\{(\R_{m_1},1+t_{111}z^{-m_1})\})$ 时，但该图中没有射线。因此特别地，上述论证表明 $\log f_{\mathfrak{d}}$ 模 $J$ 只能出现有限多个项；又因 $\log f_{\mathfrak{d}}$ 的每个系数都涉及某个满足 $t_{ijk}^N\in J$ 的 $t_{ijk}$，$f_{\mathfrak{d}}=\exp(\log f_{\mathfrak{d}})$ 模 $J$ 只有有限多个项。于是 $f_{\mathfrak{d}}$ 的项数与 $e$ 无关地有界。最后，应用注 6.46 中给出的分解过程，我们可以把每条射线 $(\mathfrak{d},f_{\mathfrak{d}})$ 换成有限多条射线 $(\mathfrak{d},f_{\mathfrak{d}})$，其中 $f_{\mathfrak{d}}$ 形如 $1+cz^m$。这便证明了所需结果。
\end{proof}

\begin{remark}
上述命题可如下解释。散射图
\[
S(\mathfrak{D})=\bigcup_{e=1}^\infty S_{\tilde I_e}(\mathfrak{D})
\]
可看作环
\[
\hat A:= \Bbbk[M]\otimes_ k\big( \Bbbk[\{t_{ijk}\mid(i,j,k)\ne(1,1,1)\}]/J\big)[\![t_{111}]\!]
\]
上的散射图，且在该环中 $\theta_{\gamma,S(\mathfrak{D})}=\operatorname{Id}$。然而 $\hat A$ 含子环
\[
A:=\big( \Bbbk[M]\otimes_ k \Bbbk[\{t_{ijk}\}]/J\big)_{1+t_{111}z^{-m_1}},
\]
而 $\theta_{\gamma,\mathfrak{D}}$ 作为 $A$ 的自同构有意义。$S(\mathfrak{D})$ 有限这一事实意味着 $\theta_{\gamma,S(\mathfrak{D})}$ 作为 $A$ 的自同构也有意义，并且当然仍是恒等映射。这告诉我们，Kontsevich–Soibelman 引理（定理 6.38）在这一特定情形下在局部化环 $A$ 上有意义。
\end{remark}

\subsubsection{使结构相容}

现在我们可以构造相容结构了。照常从 $(B,P,\phi)$ 出发：$B$ 是带有正、简单奇点的二维整热带流形。我们将构造一列结构 $S_0,S_1,\dots$，其中 $S_k$ 相容到 $k$ 阶。

从 $S_0$ 开始：这里别无选择，只能置
\[
S_0=\{d_{p,\pm}\mid p\in\Delta\},
\]
其中 $d_{p,\pm}$ 表示结构定义所要求的、与奇点 $p$ 相伴的两条射线。容易看出 $S_0$ 相容到 $0$ 阶。我们来验证这一点。首先，对 $d_{p,\pm}$，我们可以取命题 6.21 的 (2) 所承诺的数 $N_{d_{p,\pm}}^0$ 与定义 6.22 的 (2) 中定义的区间 $I_{d_{p,\pm}}$ 的非零端点重合。因此我们可以取 $P^0=P$。于是唯一的连接点

% [原文 p.303]
是 $\cP$ 的顶点。若 $v$ 是顶点，则对任何 $\mathfrak{d}\in S_0$ 有 $f_{\mathfrak{d}}\equiv1\bmod I_{v,v,\sigma}^0$，故相容性是平凡的。

现在我们想归纳地构造 $S_k$，故假设已构造了相容到 $k-1$ 阶的 $S_{k-1}$。我们将添加新射线以得到相容到 $k$ 阶的 $S_k$。为此我们使用前两小节的结果。特别地，该论证是纯局部的：在每个连接点处，我们添加由 Kontsevich–Soibelman 论证所指定的一族射线。

对给定的连接点 $j$，取 $M=\Lambda_j$（即 $\Lambda$ 在点 $j$ 处的茎），并取 $P=P_{\phi,\sigma_j,\sigma}$，其中 $\sigma\in\cP_{\max}$ 含 $j$。令 $I=I_{\sigma_j,\sigma_j,\sigma}^k$。注意 $I$ 总是 $\mathfrak{m}$-准素的，其中 $\mathfrak{m}=P\setminus P^\times$。

如下构造散射图 $\mathfrak{D}_j$。$\mathfrak{D}_j$ 的元素与满足 $\mathfrak{d}\in S_{k-1}$、$x\in[0,N_{\mathfrak{d}}^k]$ 且 $\mathfrak{d}(x)=j$ 的对 $(\mathfrak{d},x)$ 一一对应。给定这样的对 $(\mathfrak{d},x)$，我们给它关联 $\mathfrak{D}_j$ 中形如
\[
\begin{cases}(\R_{\ge0}\mathfrak{d}'(x),1+c_{\mathfrak{d}}z^{m_{\mathfrak{d},x}}),&x=0\\ (\R\,\mathfrak{d}'(x),1+c_{\mathfrak{d}}z^{m_{\mathfrak{d},x}}),&x\ne0\end{cases}
\]
的元素。这里 $\mathfrak{d}'(x)\in\Lambda_{\mathfrak{d}(x)}=\Lambda_j$ 表示 $\mathfrak{d}$ 在 $x$ 处的切向量。现在需要考虑两种情形。

$\dim\sigma_j=2$ 且 $\dim\sigma_j=0$：此时 $\mathfrak{D}_j$ 实际上是 $\Bbbk[P]/I$ 上的散射图。除了对 $\mathfrak{d}\in\mathfrak{D}_j$ 要求 $f_{\mathfrak{d}}\equiv1\bmod\mathfrak{m}$ 之外，所有条件都由结构中射线的定义显然成立。然而，若 $\mathfrak{d}$ 对应于对 $(\bar{\mathfrak{d}},x)$ 且 $\dim\sigma_j=2$，则 $\bar{\mathfrak{d}}$ 不可能是初始射线，故由定义 6.22 的 (3) 有 $\operatorname{ord}_{\sigma_j}m_{\bar{\mathfrak{d}},x}>0$。于是 $1+c_{\bar{\mathfrak{d}}}z^{m_{\bar{\mathfrak{d}},x}}\equiv1\bmod\mathfrak{m}$。另一方面，若 $\dim\sigma_j=0$，且 $\mathfrak{d}$ 对应于 $(\bar{\mathfrak{d}},x)$，则存在含 $j$ 的 $\sigma\in\cP_{\max}$ 使 $\operatorname{ord}_\sigma m_{\bar{\mathfrak{d}},x}>0$——若 $\bar{\mathfrak{d}}$ 不是初始射线，这由定义 6.22 的 (3) 得到；若 $\bar{\mathfrak{d}}$ 是初始射线，则由引理 6.19 的 (2) 得到（注意连接点永不在 $\partial B$ 中）。所以同样有 $1+c_{\bar{\mathfrak{d}}}z^{m_{\bar{\mathfrak{d}},x}}\equiv1\bmod\mathfrak{m}$。

于是我们得到有限散射图 $S_I(\mathfrak{D}_j)$。注意定理 6.38 中给出的 $S_I(\mathfrak{D}_j)$ 的特定构造总是添加满足 $f_{\mathfrak{d}}$ 形如 $1+c_{\mathfrak{d}}z^{m_{\mathfrak{d}}}$ 的射线 $(\mathfrak{d},f_{\mathfrak{d}})$，并且我们可以假设对每条添加的射线有 $m_{\mathfrak{d}}\in I_{\sigma_j,\sigma_j,\sigma}^{k-1}$。确实，$S_{k-1}$ 相容这一事实意味着对绕原点的回路 $\gamma$ 有
\[
\theta_{\gamma,\mathfrak{D}_j}\equiv\operatorname{Id}\bmod I_{\sigma_j,\sigma_j,\sigma}^{k-1}.
\]
现在对每个 $(\mathfrak{d},1+c_{\mathfrak{d}}z^{m_{\mathfrak{d}}})\in S_I(\mathfrak{D}_j)\setminus\mathfrak{D}_j$，令 $(\bar{\mathfrak{d}},f_{\bar{\mathfrak{d}}})$ 为 $B$ 上如下得到的射线：取 $\bar{\mathfrak{d}}(0)=j$，$\bar{\mathfrak{d}}'(0)=-\bar m_{\mathfrak{d}}$，$f_{\bar{\mathfrak{d}}}=1+c_{\mathfrak{d}}z^{m_{\bar{\mathfrak{d}}}}$，其中 $m_{\bar{\mathfrak{d}}}$ 由 $m_{\mathfrak{d}}$ 平行移动得到。由命题 6.21 可知这是一条射线。置
\begin{equation}
S_j=\{(\bar{\mathfrak{d}},f_{\bar{\mathfrak{d}}})\mid\mathfrak{d}\in S_I(\mathfrak{D}_j)\setminus\mathfrak{D}_j\}.
\tag{6.10}
\end{equation}

$\dim\sigma_j=1$：若边 $\sigma_j$ 含一个奇点，则不妨假设 $j\in d_{p,+}(I_{d_{p,+}}^{\mathrm{init}})$（而非 $d_{p,-}(I_{d_{p,-}}^{\mathrm{init}})$）。则对任何含 $\sigma_j$ 的 $\sigma$ 有 $\operatorname{ord}_\sigma m_-=0$，故 $f_{d_{p,+}}\not\equiv1\bmod\mathfrak{m}$。然而这是唯一造成问题的射线。特别地，给定 $M_\R$ 中的路径 $\gamma$，$\theta_{\gamma,\mathfrak{D}_j}$ 在环 $\Bbbk[P]/I$ 中没有意义，因为 $f_{d_{p,+}}$ 不可逆。它只在环 $(\Bbbk[P]/I)_{f_{d_{p,+}}}$ 中有意义，但我们尚未证明该环上版本的定理 6.38。不过我们可以用注 6.46 与命题 6.47，以另一种方式解释 $\mathfrak{D}_j$ 使之有意义。

% [原文 p.304]
特别地，把注 6.46 的构造应用于散射图
\[
\mathfrak{D}_j=\{d\in\mathfrak{D}_j\mid d\text{ 是直线}\}.
\]
我们如该注中那样把它换为等价散射图
\[
\mathfrak{D}_j=\Big\{\Big(\R_{\bar m_i},\prod_{j=1}^{p_i}\prod_{k=1}^{\ell_{ij}}(1+c_{ijk}z^{-m_{ijk}})\Big)\Big\},
\]
其中我们可以假设 $1+c_{111}z^{-m_{111}}=1+z^{m_-}$ 就是给我们造成麻烦的项。然后我们按注 6.46 的过程得到 $S(\mathfrak{D}')$。唯一的微妙之处是：$t_{111}z^{-r(m_{111})}$ 现在在 $\phi$ 下映到 $z^{m_-}$，而它不在 $\mathfrak{m}$ 中。不过由命题 6.47 可知，若 $I\subseteq P$ 是 $\mathfrak{m}$-准素理想，则 $\phi(S(\mathfrak{D}'))$ 模 $I$ 只含有限多条射线。此外仍然有 $\theta_{\gamma,\phi(S(\mathfrak{D}'))}\equiv\operatorname{Id}\bmod I$，只不过现在是在局部化环 $(\Bbbk[P]/I)_{1+z^{m_-}}$ 中，如注 6.48 所述。

现在定义 $S_I(\mathfrak{D}_j)$ 为按注 6.46 把 $\phi(S(\mathfrak{D}'))$ 中所有射线分解所得。这是一个唯一的程序，且该散射图中所有模 $I_{\sigma_j,\sigma_j,\sigma}^{k-1}$ 非平凡的射线由归纳都已在 $\mathfrak{D}_j$ 中出现。于是如前用 (6.10) 定义 $S_j$。\hfill$\square$

令
\[
S_k=S_{k-1}\cup\bigcup_{j\in\operatorname{Joints}(S_{k-1},k)}S_j.
\]

\begin{theorem}
$S_k$ 相容到 $k$ 阶。
\end{theorem}

\begin{proof}[证明]
设 $j\in\operatorname{Joints}(S_k,k)$，并令
\[
\theta_j=\theta_{u_n,u_1}\circ\cdots\circ\theta_{u_1,u_2}:R_{\sigma_j,\sigma_j,u_1}^k\to R_{\sigma_j,\sigma_j,u_1}^k
\]
为为达到相容性必须模 $I=I_{\sigma_j,\sigma_j,\sigma_{u_1}}^k$ 为恒等映射的自同构。按假设 $\theta_j\equiv\operatorname{Id}\bmod I_{\sigma_j,\sigma_j,\sigma_{u_1}}^{k-1}$，因为所有射线 $d\in S_I(\mathfrak{D}_j)\setminus\mathfrak{D}_j$ 满足 $f_d\equiv1\bmod I_{\sigma_j,\sigma_j,\sigma_{u_1}}^{k-1}$。（若 $j\notin\operatorname{Joints}(S_k,k-1)$，则我们取 $S_I(\mathfrak{D}_j)\setminus\mathfrak{D}_j$ 为空。）此外，由 $S_I(\mathfrak{D}_j)$ 的构造，$\theta_j$ 模 $I$ 不为恒等只可能因为存在某条射线 $d\in S_k$ 使对某个 $x\ne0$ 有 $d(x)=j$，而该对 $(d,x)$ 未在 $\mathfrak{D}_j$ 的构造中出现。

该对 $(d,x)$ 未出现有两种可能的原因。第一，可能 $x\notin[0,N_d^k]$，此时 $z^{m_{d,x}}\in I_{\sigma_j,\sigma_j,\sigma_{u_1}}^k$，故 $(d,x)$ 对 $\theta_j$ 无贡献。第二，可能 $d\in S_k\setminus S_{k-1}$，此时 $z^{m_{d,x}}\in I_{\sigma_j,\sigma_j,\sigma_{u_1}}^{k-1}$。我们现在按 $\sigma_j$ 可能的维数分析。

$\dim\sigma_j=2$：此时 $I_{\sigma_j,\sigma_j,\sigma_{u_1}}^0\cdot I_{\sigma_j,\sigma_j,\sigma_{u_1}}^{k-1}\subseteq I_{\sigma_j,\sigma_j,\sigma_{u_1}}^k$。于是任何这样的射线所诱导的自同构模 $I_{\sigma_j,\sigma_j,\sigma_{u_1}}^k$ 与所有经过 $j$ 的其他射线所伴随的自同构交换。由于 $d$ 对 $\theta_j$ 贡献两次、但对应互逆的自同构，事实上 $d$ 不影响 $\theta_j\equiv\operatorname{Id}\bmod I_{\sigma_j,\sigma_j,\sigma_{u_1}}^k$。

$\dim\sigma_j=1$：此时 $d$ 或者平行于边 $\sigma_j$，或者横截于该边。在第一种情形，$d$ 所伴随的自同构与 $d_{p,\pm}$ 所伴随的自同构交换，其中 $p\in\sigma_j$ 是奇点（若存在的话），这是因为 $d$ 平行于 $d_{p,\pm}$。另一方面，设 $\sigma_\pm$ 为含 $\sigma_j$ 的极大胞腔。写 $f_{(d,x)}=1+cz^m$，可见必须有 $\operatorname{ord}_{\sigma_\pm}m=k$。

% [原文 p.305]
于是 $d$ 所伴随的自同构模 $I_{\sigma_j,\sigma_j,\sigma_{u_1}}^k$ 与对 $\theta_j$ 有贡献的任何其他自同构交换。在第二种情形，不妨设 $d$ 从极大胞腔 $\sigma_+$ 进入极大胞腔 $\sigma_-$。由构造必须有 $\operatorname{ord}_{\sigma_+}(m_{d,x})\ge k$，从而由引理 6.19 有 $\operatorname{ord}_{\sigma_-}(m_{d,x})>k$。于是该射线所伴随的自同构模 $I_{\sigma_j,\sigma_j,\sigma_{u_1}}^k$ 是平凡的，对 $\theta_j$ 没有影响。

$\dim\sigma_j=0$：此时 $d$ 从含 $j$ 的极大胞腔 $\sigma_+$ 进入含 $j$ 的胞腔 $\sigma_-$。由同样的论证，$\operatorname{ord}_{\sigma_+}(m_{d,x})\ge k$ 且 $\operatorname{ord}_{\sigma_-}(m_{d,x})>k$，故 $d$ 对 $\theta_j$ 仍无影响。
\end{proof}

\emph{定理 6.6 的证明.} 由定理 6.49，上述构造产生相容到 $k$ 阶的结构 $S_k$。由定理 6.35，这给出一个在 $O_k$ 上平坦的概形 $\check X_k(B,P)$。此外由构造有 $\check X_k(B,P)\times_{R_k}R_{k-1}\cong\check X_{k-1}(B,P)$。于是我们可以对所有 $k$ 取概形 $\check X_k(B,P)$ 的极限，得到在 $\operatorname{Spf} \Bbbk[\![t]\!]$ 上平坦的形式概形 $\check X(B,P)$，正如所愿。\hfill$\square$

\subsection{评注与推广}

前两节给出的论证与 Kontsevich 和 Soibelman 给出的原始论证有许多共同之处。不过他们不使用多面体分解，并且不是构造一族退化的簇，而是通过类似的粘合构造构造一个刚性解析空间。控制这两种粘合构造的数据非常相似。但一个关键区别是，Kontsevich 与 Soibelman 在我们会称为扇图景的框架中工作。因此射线不是直的，而必须取为某些函数关于某个特殊度量的梯度流线。于是他们的构造看起来不如这里的构造那么“热带”，这里所有射线的并看起来像热带树之并。忽略多面体分解这一离散数据也使构造更易描述，因为不必做逐情形分析。尽管如此，他们的论证有一个不惬意的方面：他们需要构造度量的选取，使得没有梯度流线回到奇点的某个小邻域。这是保证构造收敛所需要的。

这一问题在此也会带来麻烦，但都被涵盖在命题 6.47 中。虽然这里的论证简单，但它当然用到了 [45] 的全部力量。在 [49] 中，我们给出了另一种更长但更初等的论证来处理这一情形。无论如何，要避免奇点造成的某些收敛性技术问题似乎是不可能的。

[49] 一文给出了所有维数的结果。本章采取的方法可视为 Kontsevich 与 Soibelman 的版本与 [49] 的方法的混合。在 [49] 中，一般地不是考虑射线，而是考虑从奇点轨迹发出的余维一墙，其上附有决定粘合的函数。这样的对象在拓扑上可能非常复杂，因此，逐阶地用积木构造墙、用位于 $\cP$ 的胞腔内的墙代替可能有无限伸展的对象，被证明要容易得多。这已经使 [49] 读起来有些困难。在那里连接点是余维二的多面体，而在每个连接点处添加新墙。

远为微妙的是那些直接来自奇点轨迹的墙上零阶项所引起的收敛性问题。特别地，保证定理 6.38 对位于 $\cP$ 的余维一与余维二胞腔内的连接点有合适版本的论证非常困难，长达约 40 页。有希望的是，随着对高维情形有更好的计数理解，这些论证即使不变得更容易，也可能变得更具概念性。另一方面，我不知道如何把 Kontsevich 与 Soibelman 的方法推广到高维；收敛性问题似乎更难。另一个出现的问题是，不同连接点相交处的余维三多面体也会起作用。这样的多面体称为\emph{间隙}。幸运的是，算法相当容易地处理在一个间隙处相交的各个连接点之间的相互作用；定理 5.39 证明中的论证使用了非常类似的想法。

话虽如此，我们能指望这一构造完成什么？要点在于，我们有了由 $(B,P,\phi)$ 决定的簇退化的一个描述，其语言本质上是 Maslov 指标为零的热带圆盘。然后可以通过粘合这样的热带圆盘得到有理热带曲线，正如第 5 章所进行的那样。需要执行第 5 章中周期计算的一个更一般的版本，以看清这些周期计算如何提取粘合这些圆盘的所有方式。一旦完成这一点，就应当对一般的 B-模型周期计算有精确的热带解释。把它与使用目前正在发展的对数 Gromov–Witten 不变量理论、按第 4 章的精神所作的 Gromov–Witten 不变量计算相结合，人们便可望证明镜像对称。在写作之时，这个方向上仍有许多工作要做。

% [原文 p.306]
\subsection{参考文献与延伸阅读}

本节材料取自本人与 Siebert 的合作工作，§6.3.2 的材料除外，那一部分叙述的是 [45] 中完成的工作；另可参见 [44] 的综述。要了解更多 Gross–Siebert 纲领的细节，[47]、[41] 与 [50] 是作为介绍性文章而写的。

纲领本身的技术细节已在 [48]、[51] 与 [49] 中发展。其中最后一篇给出了光滑化结果，读者在读完本章之后应当可以读懂它，而无需先读 [48]。

如前所述，本章给出的论证可视为 Kontsevich 与 Soibelman 在 [70] 中在刚性解析情形下给出的原始论证的一个变体。自该文写出以来，作者们已发现最初在 [70] 中出现的自同构群（这里称为热带顶点群）在跨墙公式中的应用：见 [71]。此外，这里讨论的散射图在箭图表示的模空间的 Euler 示性数方面也有解释：见 Reineke 的 [94] 与 [95]。另可参见 Gaiotto、Moore 与 Neitzke 的相关工作 [29]、[30]。本章描述的想法最近还被用于本人与 Paul Hacking、Sean Keel 的合作（即将发表）中，以证明 Looijenga 关于尖点奇点可光滑化的猜想，并为 K3 曲面构造典范 theta 函数。

% [Tutor Note: §6.5 是本章的“参考文献与延伸阅读”小节：按全书惯例只译正文行文，文献编号保留原文方括号形式，不逐条翻译文献条目。]

% [Tutor Note: 本块环境序列为定义 6.37、定理 6.38、注 6.39、定义 6.40、例 6.41、定理 6.42、例 6.43（复数）、定理 6.44、注 6.45、注 6.46、命题 6.47、注 6.48、定理 6.49，自动编号须接 c06a–c06c 的 36 个定理类环境之后；§6.3.2 与 §6.3.3 的 \subsubsection 编号依赖前块已写 §6.3.1。]
% [Tutor Note: 全块区分两种 Fraktur 记号：$\mathfrak{D}$ 表散射图，$\mathfrak{d}$ 表射线或直线；$\R_m$ 表从原点出发、方向为 $m$ 的射线，$\R_{\ge0}m$ 表其方向锥。转录核对时勿混。]
% [Tutor Note: 原文 pp.292–306 中被大量引用的定理 6.6、6.21、6.22、6.29、6.35、5.39 等均不在本块内，属前块或第 5 章，引用按纯文本处理。]
